\documentclass[14pt,a4paper]{scrartcl} 
\usepackage[utf8]{inputenc}
\usepackage[english,russian]{babel}
\usepackage{indentfirst}
\usepackage{misccorr}
\usepackage{graphicx}
\usepackage{amsmath}
\usepackage{minted}
\usepackage{geometry}
\usepackage{float}

\geometry{left=3cm, right=1.5cm, top=2cm, bottom=3cm}

\begin{document}
\thispagestyle{empty}
\setcounter{page}{2}
\tableofcontents
\newpage

\section{Введение}
\label{sec:intro}

\subsection{Формулировка задачи}

Требуется написать программу для решения системы линейных алгебраических уравнений методом Крамера. Количество неизвестных в системе может быть 2 или 3 (вариант 4). Программа должна проверять корректность входных данных и обрабатывать случаи, когда метод не применим.

Метод Крамера применяется к системе вида:
\begin{equation}
\begin{cases}
a_{11}x_1 + a_{12}x_2 + \dots + a_{1n}x_n = b_1 \\
a_{21}x_1 + a_{22}x_2 + \dots + a_{2n}x_n = b_2 \\
\vdots \\
a_{n1}x_1 + a_{n2}x_2 + \dots + a_{nn}x_n = b_n
\end{cases}
\end{equation}

где число уравнений равно числу неизвестных.

\subsection{Математическая теория}

Если главный определитель матрицы коэффициентов $\Delta \neq 0$, то решение единственно и находится по формулам:

\begin{equation}
x_i = \frac{\Delta_i}{\Delta}, \quad i = 1, 2, \dots, n
\end{equation}

где $\Delta$ --- определитель матрицы коэффициентов, $\Delta_i$ --- определитель матрицы, полученной заменой $i$-го столбца на столбец свободных членов.

Для матрицы второго порядка:
\begin{equation}
\Delta = \begin{vmatrix} a_{11} & a_{12} \\ a_{21} & a_{22} \end{vmatrix} = a_{11}a_{22} - a_{12}a_{21}
\end{equation}

Для матрицы третьего порядка используется разложение по первой строке:
\begin{equation}
\Delta = a_{11}M_{11} - a_{12}M_{12} + a_{13}M_{13}
\end{equation}

где $M_{1k}$ --- миноры соответствующих элементов.

\section{Ход работы}
\label{sec:exp}

\subsection{Описание алгоритма программы}

Программа работает по следующему алгоритму:

\begin{enumerate}
\item Запрашивает количество неизвестных (2 или 3) с проверкой корректности.
\item Считывает матрицу коэффициентов построчно, проверяя что введены числа.
\item Считывает вектор свободных членов с проверкой.
\item Выводит введённую систему для проверки пользователем.
\item Вычисляет главный определитель $\Delta$ рекурсивным методом.
\item Если $\Delta = 0$, сообщает об ошибке и завершает работу.
\item Если $\Delta \neq 0$, вычисляет для каждой переменной определитель вспомогательной матрицы и применяет формулу Крамера.
\item Выводит решение на экран.
\end{enumerate}

\subsection{Код приложения}
\label{sec:exp:code}

\begin{minted}[fontsize=\small, breaklines]{python}
def read_size():
    while True:
        raw = input("введите количество неизвестных (2 / 3): ").strip()
        if raw not in ("2", "3"):
            print("ошибка! метод Крамера работает только для 2 или 3 неизвестных!")
            continue
        return int(raw)


def read_float(prompt):
    while True:
        raw = input(prompt).strip()
        try:
            return float(raw)
        except ValueError:
            print("ошибка! введите число!")


def read_matrix(n):
    print("\nвведите коэффициенты матрицы по строкам")
    matrix = []
    for i in range(n):
        row = []
        for j in range(n):
            value = read_float(f"  a[{i + 1}][{j + 1}] = ")
            row.append(value)
        matrix.append(row)
    return matrix


def read_vector(n):
    print("\nвведите свободные члены")
    vector = []
    for i in range(n):
        value = read_float(f"  b[{i + 1}] = ")
        vector.append(value)
    return vector


def determinant(matrix):
    n = len(matrix)
    if n == 1:
        return matrix[0][0]
    if n == 2:
        return matrix[0][0] * matrix[1][1] - matrix[0][1] * matrix[1][0]
    total = 0.0
    for col in range(n):
        minor = []
        for row in range(1, n):
            minor_row = matrix[row][:col] + matrix[row][col + 1:]
            minor.append(minor_row)
        sign = 1 if col % 2 == 0 else -1
        total += sign * matrix[0][col] * determinant(minor)
    return total


def replace_column(matrix, column_index, new_column):
    new_matrix = []
    for row in range(len(matrix)):
        new_row = matrix[row][:]
        new_row[column_index] = new_column[row]
        new_matrix.append(new_row)
    return new_matrix


def cramer_solve(matrix, free_terms):
    main_det = determinant(matrix)
    if abs(main_det) < 1e-12:
        return None, main_det
    solution = []
    for i in range(len(matrix)):
        modified = replace_column(matrix, i, free_terms)
        det_i = determinant(modified)
        solution.append(det_i / main_det)
    return solution, main_det


def print_system(matrix, free_terms):
    n = len(matrix)
    print("\nвведённая система уравнений:")
    for i in range(n):
        terms = []
        for j in range(n):
            terms.append(f"{matrix[i][j]:.3f}*x{j + 1}")
        line = " + ".join(terms)
        print(f"  {line} = {free_terms[i]:.3f}")


def main():
    print("Решение СЛУ методом Крамера")
    print("-" * 40)
    n = read_size()
    matrix = read_matrix(n)
    free_terms = read_vector(n)
    print_system(matrix, free_terms)
    solution, main_det = cramer_solve(matrix, free_terms)
    print(f"\nглавный определитель: {main_det:.6f}")
    if solution is None:
        print("система не имеет единственного решения (главный определитель = 0)")
        print("метод Крамера не применим!")
        return
    print("\nрешение системы:")
    for i, value in enumerate(solution):
        print(f"  x{i + 1} = {value:.6f}")


if __name__ == "__main__":
    main()
\end{minted}

\subsection{Формула метода Крамера}
\label{sec:mathexample}

Формула для нахождения каждой переменной:
\begin{equation}\label{eq:cramer}
x_i = \frac{\Delta_i}{\Delta}, \quad \Delta \neq 0
\end{equation}

Ссылка на формулу~\eqref{eq:cramer}.

\section{Примеры работы программы}
\label{sec:examples}

\subsection{Пример 1: Система из двух уравнений}

Рассмотрим систему:
\begin{equation}
\begin{cases}
2x_1 + 3x_2 = 8 \\
x_1 - x_2 = 1
\end{cases}
\end{equation}

Главный определитель: $\Delta = 2 \cdot (-1) - 3 \cdot 1 = -5 \neq 0$.
Результат: $x_1 = 2.2$, $x_2 = 1.2$.


\begin{figure}[H]
    \centering
    \includegraphics[width=0.8\textwidth]{1.png}
    \caption{Результат выполнения Примера 1}
\end{figure}

\subsection{Пример 2: Система из трёх уравнений}

Рассмотрим систему:
\begin{equation}
\begin{cases}
x_1 + 2x_2 + 3x_3 = 6 \\
4x_1 + 5x_2 + 6x_3 = 15 \\
7x_1 + 8x_2 + 10x_3 = 25
\end{cases}
\end{equation}

Главный определитель: $\Delta = -3 \neq 0$.
Результат: $x_1 = 1$, $x_2 = 1$, $x_3 = 1$.

\begin{figure}[H]
    \centering
    \includegraphics[width=0.8\textwidth]{2.png}
    \caption{Результат выполнения Примера 2}
\end{figure}

\subsection{Пример 3: Вырожденная система}

Рассмотрим систему:
\begin{equation}
\begin{cases}
x_1 + 2x_2 = 5 \\
2x_1 + 4x_2 = 9
\end{cases}
\end{equation}

Главный определитель: $\Delta = 0$. Метод Крамера неприменим.

\begin{figure}[H]
    \centering
    \includegraphics[width=0.8\textwidth]{3.png}
    \caption{Результат выполнения Примера 3}
\end{figure}

\section{Заключение}

В ходе работы реализована программа на языке Python, которая решает системы линейных алгебраических уравнений методом Крамера для систем из двух и трёх неизвестных. Программа корректно проверяет входные данные, вычисляет определители и обрабатывает случай вырожденной системы.

\addcontentsline{toc}{section}{Список литературы}
\begin{thebibliography}{9}
\bibitem{Korn-1984}Корн Г., Корн Т. Справочник по математике. \newblock --- Москва: Наука, 1984. 831~с.
\bibitem{Faddeev-1963}Фаддеев Д.К., Фаддеева В.Н. Вычислительные методы линейной алгебры. \newblock --- Москва: Физматгиз, 1963. 734~с.
\end{thebibliography}

\end{document}